\subsection*{Why the restricted interface works: expressiveness vs performance}

Writing every load and address calculation separately gives a programmer many choices, but it also leaves the compiler to recover which operations form one tile transfer. Croqtile asks the program to describe that region explicitly. The restriction is useful when the desired movement fits a regular shape, because its geometry is available before individual transfers are generated.

Croqtile makes tile geometry and layout explicit through structured movement operations. This gives the compiler the information it needs to lower a transfer efficiently and avoids repeating per-element address choreography.

\begin{figure}[htbp]\centering
\includegraphics[width=\linewidth,height=0.56\textheight,keepaspectratio]{f09-fig8_expressiveness_light.png}
\caption{Structured tile operations expose geometry and layout for efficient lowering; performance still depends on the chosen tiles and execution configuration.}
\end{figure}

The trade is explicit: structured operations constrain how a transfer is described. They support efficient lowering, while actual performance still depends on tile sizes, alignment, layouts, and synchronization.
